\documentclass[11pt,a4paper]{report}

\usepackage[utf8]{inputenc}
\usepackage[T1]{fontenc}
\usepackage{lmodern}
\usepackage{graphicx}
\usepackage{xcolor}
\usepackage{hyperref}
\usepackage{geometry}
\usepackage{booktabs}
\usepackage{listings}
\usepackage{tcolorbox}
\usepackage{tikz}
\usetikzlibrary{shapes.geometric, arrows, positioning, shadows}

\geometry{margin=1in}

\definecolor{primary}{RGB}{0, 102, 204}
\definecolor{secondary}{RGB}{102, 102, 102}
\definecolor{codebg}{RGB}{245, 245, 245}

\hypersetup{
    colorlinks=true,
    linkcolor=primary,
    filecolor=magenta,      
    urlcolor=primary,
    pdftitle={XDP-L2-Guard Architecture Overview},
    pdfpagemode=FullScreen,
}

\lstset{
    backgroundcolor=\color{codebg},
    basicstyle=\ttfamily\small,
    breaklines=true,
    frame=single,
    rulecolor=\color{secondary},
    keywordstyle=\color{primary},
}

\title{\textbf{XDP-L2-Guard Architecture Overview}}
\author{Development Team}
\date{\today}

\begin{document}

\maketitle
\tableofcontents

\chapter{Introduction}
XDP-L2-Guard follows a decoupled architecture, separating the high-performance data plane (Kernel Space) from the management and configuration logic (User Space).

\section{System Components}

\begin{center}
\begin{tikzpicture}[
    node distance=1.5cm,
    block/.style={rectangle, draw, fill=blue!10, text width=5em, text centered, rounded corners, minimum height=3em, drop shadow},
    map/.style={cylinder, draw, fill=yellow!10, text width=4em, text centered, shape border rotate=90, minimum height=3em, drop shadow},
    cloud/.style={draw, ellipse, fill=green!10, text width=5em, text centered, minimum height=3em, drop shadow},
    arrow/.style={thick,->,>=stealth}
]

% User Space
\node [block] (cli) {guard.py / guard.sh};
\node [block, below of=cli] (loader) {Loader / Controller};
\node [map, right=2cm of loader] (maps) {BPF Maps};

\node[draw, dashed, inner sep=0.5cm, label=above:User Space, fit=(cli) (loader) (maps)] (userspace) {};

% Kernel Space
\node [block, below=2cm of loader] (filter) {XDP Filter Program};
\node [block, left=1cm of filter] (nic) {NIC Driver};

\node [cloud, below left=1cm and 0.5cm of filter] (discard) {Discard};
\node [cloud, below=1cm of filter] (stack) {Linux Stack};
\node [cloud, below right=1cm and 0.5cm of filter] (bounce) {Bounce};
\node [cloud, right=1cm of filter] (redirect) {Redirect};

\node[draw, dashed, inner sep=0.5cm, label=below:Kernel Space, fit=(filter) (nic) (discard) (stack) (bounce) (redirect)] (kernel) {};

% Connections
\draw [arrow] (cli) -- node[right] {\scriptsize CLI} (loader);
\draw [arrow] (loader) -- node[above] {\scriptsize bpf\_map\_update} (maps);
\draw [arrow] (maps) -- node[right] {\scriptsize lookup} (filter);
\draw [arrow] (nic) -- node[above] {\scriptsize DMA} (filter);

\draw [arrow] (filter) -- node[left] {\scriptsize XDP\_DROP} (discard);
\draw [arrow] (filter) -- node[right] {\scriptsize XDP\_PASS} (stack);
\draw [arrow] (filter) -- node[right] {\scriptsize XDP\_TX} (bounce);
\draw [arrow] (filter) -- node[above] {\scriptsize XDP\_REDIRECT} (redirect);

\end{tikzpicture}
\end{center}

\chapter{Data Plane (eBPF/XDP)}

The data plane is implemented in C and compiled into eBPF bytecode. It hooks into the network driver's NAPI loop using XDP (eXpress Data Path).

\section{Key Features}
\begin{itemize}
    \item \textbf{Direct Access}: Operates on \texttt{xdp\_md} context, providing direct access to packet data.
    \item \textbf{Stateless}: The filter is primarily stateless, relying on BPF maps for rule lookup.
    \item \textbf{Performance}: Bypasses \texttt{sk\_buff} allocation, leading to significant performance gains under high traffic.
\end{itemize}

\section{Hook Points}
\begin{itemize}
    \item \texttt{SEC("xdp")}: The main entry point \texttt{xdp\_drop\_logic}.
\end{itemize}

\chapter{Control Plane (Python/Bash)}

The control plane manages the lifecycle of the eBPF program and provides a user-friendly interface for rule management.

\section{Components}
\begin{itemize}
    \item \textbf{Loader (\texttt{src/control\_plane/loader.py})}: Loads the eBPF object file into the kernel and attaches it to the specified interface. It also populates the initial maps.
    \item \textbf{CLI (\texttt{src/control\_plane/guard.py})}: An \texttt{iptables}-like interface to interact with the \texttt{action\_map}.
    \item \textbf{Bash Wrapper (\texttt{scripts/guard.sh})}: A lightweight version using \texttt{bpftool} for environments where Python might not be preferred.
\end{itemize}

\chapter{BPF Maps}

Maps are the bridge between User Space and Kernel Space.

\begin{enumerate}
    \item \textbf{\texttt{action\_map} (Hash Map)}:
    \begin{itemize}
        \item \textbf{Key}: Destination IP (\texttt{\_\_u32}).
        \item \textbf{Value}: \texttt{struct action\_cfg} containing:
        \begin{itemize}
            \item \texttt{target}: Action to perform (DROP, PASS, TX, REDIRECT, NAT).
            \item \texttt{new\_ip}: Destination IP for NAT.
            \item \texttt{ifindex}: Target interface index for REDIRECT.
            \item \texttt{dropped\_packets}: Counter for monitoring.
        \end{itemize}
    \end{itemize}
    \item \textbf{\texttt{dev\_map} (Device Map)}:
    \begin{itemize}
        \item Used for \texttt{XDP\_REDIRECT} to efficiently forward packets to other interfaces.
    \end{itemize}
    \item \textbf{\texttt{total\_dropped} (Array Map)}:
    \begin{itemize}
        \item A single-element array to track total dropped packets across all rules.
    \end{itemize}
\end{enumerate}

\chapter{Packet Flow Logic}

\begin{enumerate}
    \item Packet arrives at the NIC.
    \item XDP program is triggered before \texttt{sk\_buff} is created.
    \item Program parses Ethernet and IP headers.
    \item Program lookups the Destination IP in \texttt{action\_map}.
    \item If a match is found, the specified action is performed.
    \item If no match, \texttt{XDP\_PASS} is returned, letting the packet continue to the Linux network stack.
\end{enumerate}

\end{document}
